\documentclass{article}

\usepackage[english]{babel}
\usepackage{enumerate}
\usepackage[letterpaper,top=2cm,bottom=2cm,left=3cm,right=3cm,marginparwidth=1.75cm]{geometry}
\usepackage{amsmath}
\usepackage{graphicx}
\usepackage[colorlinks=true, allcolors=blue]{hyperref}
\usepackage{titling}
\usepackage{listings}
\usepackage{xcolor}
\usepackage{amssymb}
\usepackage{float}
\usepackage{booktabs}
\usepackage{longtable}

\definecolor{codegreen}{rgb}{0,0.6,0}
\definecolor{codegray}{rgb}{0.5,0.5,0.5}
\definecolor{codepurple}{rgb}{0.58,0,0.82}
\definecolor{backcolour}{rgb}{0.95,0.95,0.92}

\lstdefinestyle{mystyle}{
    backgroundcolor=\color{backcolour},
    commentstyle=\color{codegreen},
    keywordstyle=\color{magenta},
    numberstyle=\tiny\color{codegray},
    stringstyle=\color{codepurple},
    basicstyle=\ttfamily\footnotesize,
    breakatwhitespace=false,
    breaklines=true,
    captionpos=b,
    keepspaces=true,
    numbers=left,
    numbersep=5pt,
    showspaces=false,
    showstringspaces=false,
    showtabs=false,
    tabsize=2
}

\lstset{style=mystyle}

\setlength{\parskip}{\baselineskip}%
\setlength{\parindent}{0pt}%

\begin{document}

\title{Autonomous Forest Drone\\[0.4em]\large Developer Handover Document}
\date{April 2026}
\author{Elina Rosato, Danielis Maizelis\\[0.3em]
\small Kristianstad University, Sweden}

\maketitle
\newpage

\tableofcontents
\newpage

% =============================================================
\section{Introduction}
% =============================================================

This document provides the information necessary for a developer or researcher to take over the Autonomous Forest Drone project and continue its development. It covers the full system from hardware assembly to software setup and progressive commissioning, and should be read alongside the companion \textbf{User Manual}, which covers day-to-day operation, pre-flight procedures, and troubleshooting.

\textbf{Scope of this document:}
\begin{itemize}
    \item Physical assembly of the drone from parts
    \item All wiring connections with connector types on both ends
    \item Pixhawk firmware flashing, calibration, and parameter configuration
    \item Software environment setup on the Jetson
    \item Repository setup and architecture
    \item Progressive commissioning and testing sequence
    \item Known limitations and suggested future work
\end{itemize}


\textbf{GitHub organisation:} \url{https://github.com/autonomous-forest-drone}

% =============================================================
\section{System Architecture}
% =============================================================

The system consists of three layers:

\begin{enumerate}
    \item \textbf{Flight platform}: Holybro S500 V2 quadcopter frame with Pixhawk 6C running PX4 firmware. The Pixhawk handles all low-level flight control, sensor fusion, and safety features.
    \item \textbf{Companion computer}: NVIDIA Jetson Orin Nano 8\,GB running ROS~2 Humble with MAVROS. The Jetson executes autonomous mission scripts that communicate with the Pixhawk over MAVLink via a UART serial link (Jetson GPIO TX/RX pins to Pixhawk TELEM2).
    \item \textbf{Ground station}: A laptop running QGroundControl for monitoring and parameter management, and the \texttt{ground-control-station} repository scripts for receiving telemetry from the drone.
\end{enumerate}

% =============================================================
\section{Hardware Assembly}
% =============================================================

\subsection{Bill of Materials}

All components required to build and operate the drone are listed below.

\subsubsection{Purchased Components}

\begin{longtable}{p{4.5cm}p{4.5cm}p{4cm}}
\toprule
\textbf{Component} & \textbf{Part Name / Model} & \textbf{Notes} \\
\midrule
\endhead
Frame kit & Holybro S500 V2 Quadcopter Kit & Includes frame, arms, landing gear, PDB, motors, ESCs, and hardware \\
Flight controller & Holybro Pixhawk 6C & \\
GPS module & Holybro M10 GPS with compass & \\
Power module & Holybro PM02 & Voltage and current sensing \\
Telemetry radio (air) & RFD900x or SiK Radio & 915\,MHz or 433\,MHz \\
Telemetry radio (ground) & Matching RFD900x or SiK unit & Must match frequency band \\
RC receiver & FlySky FS-iA6B & PPM mode \\
RC transmitter & FlySky FS-i6 & 6-channel, 2.4\,GHz AFHDS \\
Companion computer & NVIDIA Jetson Orin Nano 8\,GB Developer Kit & Includes case \\
Camera & Arducam IMX219 & 8\,MP, 62$^\circ$ FOV, MIPI CSI, fixed focus \\
Main battery & 4S LiPo 14.8\,V 5000\,mAh & XT60 connector \\
Battery charger & 4S LiPo balance charger & Balance charging required \\
XT60 Y-splitter cable & XT60 female to 2$\times$ XT60 male & Splits battery to power module and Jetson supply \\
Jetson power adapter cable & XT60 male to 5.5$\times$2.5\,mm barrel jack & Direct connection from Y-splitter to Jetson Orin Nano power input; no converter needed (Jetson accepts 9--20\,V) \\
Vibration dampening plate & M3-compatible foam or rubber standoff plate & Between Pixhawk and Jetson case \\
CSI ribbon cable & 22-pin (Jetson) to 15-pin (camera), 0.5\,mm pitch & Arducam IMX219 to Jetson CSI0 (camera0) \\
Jetson--Pixhawk UART cable & Female Dupont 3-pin (Jetson GPIO) to JST-GH 6-pin (Pixhawk TELEM2) & Custom cable; TX, RX, GND \\
Motor mounting screws (extended) & M3 $\times$ (stock length + 6\,mm) & See Section~\ref{sec:propguard}; do not exceed this length \\
Velcro straps ($\times$2) & Heavy-duty adhesive velcro & Battery and Jetson mounting \\
Zip ties & Assorted & Cable management \\
\bottomrule
\end{longtable}

\subsubsection{3D Printed Parts}

The following parts were designed and 3D printed for this build. Print files are stored in the project GitHub organisation.

\begin{tabular}{p{4cm}p{3.5cm}p{6cm}}
\toprule
\textbf{Part} & \textbf{Quantity} & \textbf{Notes} \\
\midrule
Propeller guard & 4 & One per motor arm \\
Camera mount & 1 & Attaches to front of frame; holds Arducam IMX219 facing forward \\
GPS mast & 1 & Custom height required -- the original Holybro mast is too short when the Jetson is installed. The mast must raise the GPS module high enough to clear the Jetson and its case to avoid magnetic interference. \\
\bottomrule
\end{tabular}

\subsection{Tools Required}

\begin{itemize}
    \item Hex key set (M2, M2.5, M3)
    \item Philips and flathead screwdrivers
    \item Soldering iron and solder (for any power connector modifications)
    \item Digital multimeter (continuity and voltage checks)
    \item LiPo voltage checker (balance connector)
    \item 3D printer (for the custom parts listed above)
    \item Computer with QGroundControl installed
\end{itemize}

\subsection{Assembly Guide}

Assemble the drone in the order described below. Do not connect the main battery until Section~\ref{sec:motortest}.

\subsubsection{Step 1 -- Frame}

Assemble the Holybro S500 V2 frame following the kit instructions:
\begin{enumerate}
    \item Attach the four arms to the central frame body.
    \item Attach the landing gear legs to the underside.
    \item Mount the Power Distribution Board (PDB) in the centre of the frame. The PDB is included in the S500 kit.
\end{enumerate}

\subsubsection{Step 2 -- Motors}

\begin{enumerate}
    \item Install the four Holybro 2216 880KV brushless motors on the arm motor mounts following the CW/CCW positions specified for the S500 V2 kit (see the motor diagram in the User Manual, Appendix~B).
    \item \textbf{Do not fully tighten the motor screws yet} -- the propeller guards are installed in the next step and require replacing the screws.
\end{enumerate}

\subsubsection{Step 3 -- Propeller Guards}
\label{sec:propguard}

\begin{enumerate}
    \item 3D print four propeller guards.
    \item Place one guard around each motor mount before the final motor screw installation.
    \item The guards add thickness between the motor and the arm. The standard M3 motor screws included in the S500 kit are too short to grip the motor securely with the guard in place. \textbf{Replace the motor screws with M3 screws that are 6\,mm longer than the stock screws.}
    \item \textbf{Important:} Do not use screws longer than necessary. If the screw protrudes too far into the motor casing, it will damage the motor windings. 6\,mm of additional length is the correct amount for the printed guards used in this build.
    \item Tighten all motor screws with a hex key.
\end{enumerate}

\subsubsection{Step 4 -- Battery Mount and Power Module}

\begin{enumerate}
    \item Attach adhesive velcro to the battery holder on the underside of the frame and to the battery.
    \item Mount the Holybro PM02 Power Module between the battery output and the PDB. The power module has an XT60 female input (battery side) and an XT60 male output (PDB side). It also has a 6-pin JST-GH cable that connects to the Pixhawk POWER1 port for voltage and current sensing.
    \item The battery connects to the power module via the XT60 Y-splitter (see Section~\ref{sec:wiring}).
\end{enumerate}

\subsubsection{Step 5 -- Jetson Orin Nano}

\begin{enumerate}
    \item Attach heavy-duty adhesive velcro to the top of the central frame body and to the bottom of the Jetson Orin Nano case.
    \item Seat the Jetson (in its case) on the velcro on top of the frame. Centre it for weight balance.
    \item Route the Jetson power adapter cable (XT60 to barrel jack) to the rear of the frame where the XT60 Y-splitter will be accessible. The Jetson Orin Nano Developer Kit accepts 9--20\,V directly, so the LiPo voltage (14.8\,V nominal) is within spec with no converter needed.
\end{enumerate}

\subsubsection{Step 6 -- Pixhawk with Vibration Dampening}

\begin{enumerate}
    \item Attach the vibration dampening plate to the top of the Jetson case using the adhesive on the plate.
    \item Mount the Pixhawk 6C on top of the vibration dampening plate. The arrow on the Pixhawk must point toward the front of the drone.
    \item Secure the Pixhawk to the plate using double sided tape.
\end{enumerate}

\subsubsection{Step 7 -- Telemetry Radio}

Mount the air-side telemetry radio (RFD900x or SiK) on the frame using velcro or a tie. Position the antenna so it points vertically downward when the drone is in flight for optimal ground link. Route the cable to the Pixhawk TELEM1 port.

\subsubsection{Step 8 -- RC Receiver}

Mount the FlySky FS-iA6B RC receiver on the frame with its antenna wires routed away from the ESC power cables to minimise RF interference. Route the PPM signal cable to the Pixhawk RC IN port.

\subsubsection{Step 9 -- Camera Mount and Camera}

\begin{enumerate}
    \item 3D print the camera mount.
    \item Install the Arducam IMX219 camera in the mount with the lens facing forward.
    \item Attach the mount to the front of the frame.
    \item Route the 22-pin to 15-pin CSI ribbon cable from the camera toward the Jetson CSI0 connector (camera0). Take care not to kink or crease the ribbon cable; leave a gentle loop for vibration relief.
\end{enumerate}

\subsubsection{Step 10 -- RC Receiver Binding}

Bind the FlySky FS-iA6B receiver to the FlySky FS-i6 transmitter before permanently routing cables:
\begin{enumerate}
    \item Power the receiver using a temporary 5\,V source (e.g.\ connect it to the Pixhawk RC IN port and power the Pixhawk via USB from a computer).
    \item Hold the bind button on the receiver while powering it on. The LED will flash rapidly, indicating bind mode.
    \item On the FS-i6 transmitter, navigate to \textbf{System $\rightarrow$ RX Bind} and confirm. The transmitter emits a bind signal.
    \item Once bound, the receiver LED turns solid. Power cycle the receiver to confirm the bind is stored.
    \item Confirm the receiver is set to \textbf{PPM mode}: on power-on, the LED should flash twice. One flash indicates PWM mode -- switch to PPM by pressing the bind button briefly after power-on (consult the FS-iA6B manual for the exact procedure for your firmware version).
\end{enumerate}

\subsubsection{Step 11 -- GPS Mast and GPS Module}

\begin{enumerate}
    \item 3D print the custom GPS mast. The mast must be tall enough to position the GPS module above and clear of the Jetson Orin Nano. The original Holybro mast was found to be too short; the Jetson creates magnetic interference that corrupts the compass if the GPS is too close to it.
    \item Attach the mast to the rear-top of the frame.
    \item Mount the Holybro M10 GPS module on top of the mast with the arrow pointing forward (same as the Pixhawk orientation).
    \item Route the GPS cable down the mast and to the Pixhawk GPS1 port.
\end{enumerate}

% =============================================================
\section{Wiring Guide}
% =============================================================
\label{sec:wiring}

Complete all wiring after the physical assembly is done. Refer to the connection diagram in the User Manual for a visual overview. The table below lists every cable with the connector type on both ends.

\begin{longtable}{p{3.5cm}p{3.5cm}p{3.5cm}p{2.5cm}}
\toprule
\textbf{Cable} & \textbf{From} & \textbf{To} & \textbf{Connector} \\
\midrule
\endhead
Battery to Y-splitter & 4S LiPo (female XT60) & XT60 Y-splitter input & XT60 \\
Y-splitter to power module & XT60 Y-splitter output 1 & PM02 Power Module input & XT60 \\
Y-splitter to Jetson & XT60 Y-splitter output 2 & Jetson power adapter (XT60 male end) & XT60 \\
Jetson power & Jetson power adapter (barrel end) & Jetson Orin Nano power jack & 5.5$\times$2.5\,mm barrel \\
Power module to Pixhawk & PM02 telemetry output & Pixhawk 6C POWER1 & 6-pin JST-GH \\
Power module to PDB & PM02 output & PDB XT60 input & XT60 \\
ESC 1--4 signal wires & Each ESC signal lead (female Dupont/servo) & PWM signal distribution board channels 1--4 & Servo 3-pin (signal/+/GND) \\
PWM board to Pixhawk & PWM distribution board output & Pixhawk 6C MAIN OUT 1--4 & Servo rail / JST \\
GPS & Holybro M10 GPS & Pixhawk 6C GPS1 & 8-pin JST-GH \\
Telemetry radio (air) & RFD900x / SiK air unit & Pixhawk 6C TELEM1 & 6-pin JST-GH \\
RC receiver & FlySky FS-iA6B PPM output & Pixhawk 6C RC IN & 3-pin servo (PPM) \\
Jetson to Pixhawk (MAVLink) & Jetson GPIO TX/RX/GND (Dupont female) & Pixhawk 6C TELEM2 (JST-GH 6-pin) & Custom UART cable \\
Camera CSI & Arducam IMX219 & Jetson Orin Nano CSI0 (camera0) & 22-pin to 15-pin ribbon \\
\bottomrule
\end{longtable}

\textbf{Notes on the ESC wiring:}

The S500 kit ESCs each have a three-wire servo signal lead (signal, +5\,V BEC, GND). These were connected to a PWM signal distribution board (a small PCB with male pin headers for each channel). Channels 1--4 of that board receive the four ESC signal leads, and the board's output connector plugs into the Pixhawk MAIN OUT rail. This avoids routing four individual servo wires directly to the Pixhawk.

Motor-to-ESC power wiring (the thick 3-phase cables) is part of the S500 kit and soldered from the factory; verify continuity if any cable has been disturbed.

\textbf{Notes on the Jetson--Pixhawk UART link:}

The Jetson connects to the Pixhawk via a custom UART cable: Jetson GPIO TX and RX pins (40-pin header) to Pixhawk TELEM2 (JST-GH 6-pin). Wire TX on the Jetson to RX on the Pixhawk, RX on the Jetson to TX on the Pixhawk, and GND to GND. On the Jetson Orin Nano, the UART appears as \texttt{/dev/ttyTHS1}. The PX4 TELEM2 baud rate must be set to 115200 via the \texttt{SER\_TEL2\_BAUD} parameter (see Section~\ref{sec:params}). If the MAVROS connection fails on startup, verify the device node exists with \texttt{ls /dev/ttyTHS*} and that the user has serial port permissions (\texttt{sudo usermod -aG dialout \$USER}, then log out and back in).

\textbf{Cable management:}

After verifying all connections, use zip ties to bundle cables neatly along the frame arms and body. Keep high-current power cables (battery, ESC, motor) separated from signal cables (GPS, RC, telemetry, CSI) to minimise interference.

% =============================================================
\section{Pixhawk Configuration}
% =============================================================

\subsection{Firmware and Airframe}

\begin{enumerate}
    \item Download and install QGroundControl from \url{https://qgroundcontrol.com}.
    \item Connect the Pixhawk to a computer via the USB-C cable. \textbf{Do not connect the main battery.}
    \item In QGC, go to \textbf{Vehicle Setup $\rightarrow$ Firmware}. Select \textbf{PX4 Flight Stack -- Stable Release} and confirm the flash. Wait for the process to complete and the Pixhawk to reboot.
    \item After flashing, QGC prompts for airframe selection. Choose \textbf{Holybro S500} (under the Quadrotor X category). Click \textbf{Apply and Restart}.
\end{enumerate}

\textbf{Important:} Flashing firmware resets all parameters to defaults. Always run the full calibration sequence and re-apply the custom parameters after any reflash.

\subsection{Battery Configuration}

In QGC, go to \textbf{Vehicle Setup $\rightarrow$ Power} and set the battery fields to match the 4S 5000\,mAh LiPo:

\begin{tabular}{ll}
\toprule
\textbf{Field} & \textbf{Value} \\
\midrule
Number of cells & 4 \\
Full voltage (per cell) & 4.20\,V \\
Empty voltage (per cell) & 3.50\,V \\
\bottomrule
\end{tabular}

\subsection{Calibration Sequence}

Run the calibrations in the order listed below via \textbf{Vehicle Setup $\rightarrow$ Sensors}.

\begin{enumerate}
    \item \textbf{Level Horizon} -- Place the drone on a level surface and click \textbf{Level Horizon}.
    \item \textbf{Accelerometer (6-Position)} -- Follow the on-screen prompts. Place the drone in each of the six orientations indicated (level, nose down, nose up, right side down, left side down, upside down) and confirm each position.
    \item \textbf{Gyroscope} -- Place the drone on a still, level surface. Do not move it during calibration.
    \item \textbf{Compass} -- Perform outdoors, at least 2\,m from metal objects, buildings, and electronics. Follow the on-screen sphere visualisation, rotating the drone through all axes until all coloured segments are filled.
    \item \textbf{ESC Calibration} -- Remove all propellers. Go to \textbf{Vehicle Setup $\rightarrow$ Power $\rightarrow$ ESC Calibration}. Follow the QGC wizard: set the RC throttle to maximum, connect the main battery, wait for confirmation beeps from the ESCs, then lower throttle to minimum for a second set of beeps. Disconnect the battery when complete.
    \item \textbf{RC Radio} -- Go to \textbf{Vehicle Setup $\rightarrow$ Radio}. Move all sticks to their full extents and all switches through all positions. Click \textbf{Next} to save.
    \item \textbf{Flight Modes} -- Go to \textbf{Vehicle Setup $\rightarrow$ Flight Modes}. Assign Channel 5 (3-position switch) to cycle through Stabilise, Altitude Hold, and Position modes. This switch is also used to cancel Offboard mode during autonomous flight.
\end{enumerate}

\subsection{Failsafe Configuration}

In QGC, go to \textbf{Vehicle Setup $\rightarrow$ Safety} and configure:

\begin{tabular}{ll}
\toprule
\textbf{Failsafe} & \textbf{Setting} \\
\midrule
RC loss action & Return to Launch (RTL) after 1\,s \\
Low battery action & Return to Launch at 20\%; Land at 10\% \\
\bottomrule
\end{tabular}

\subsection{Parameter Configuration}
\label{sec:params}

After calibration, apply the project-specific parameters. A parameter file (\texttt{.params} export) is stored in the \texttt{drone-brain} repository (\url{https://github.com/autonomous-forest-drone}).

To apply parameters from the file:
\begin{enumerate}
    \item In QGC, go to \textbf{Vehicle Setup $\rightarrow$ Parameters}.
    \item Click the \textbf{Tools} menu (top right) and select \textbf{Load from file}.
    \item Select the \texttt{.params} file from the repository.
    \item Click \textbf{Save} and reboot the Pixhawk when prompted.
\end{enumerate}

% =============================================================
\section{Software Setup}
% =============================================================

\subsection{Jetson Orin Nano Environment}

The Jetson Orin Nano Developer Kit ships with JetPack (Ubuntu 22.04 + CUDA). Before setting up the project software, verify the OS is up to date:

\begin{lstlisting}[language=bash]
sudo apt update && sudo apt upgrade -y
\end{lstlisting}

\subsection{ROS 2 Humble and MAVROS}

ROS~2 Humble must be installed on the Jetson and is the communication backbone between the mission scripts and the Pixhawk.

\textbf{Install ROS~2 Humble} following the official instructions at \url{https://docs.ros.org/en/humble/Installation/Ubuntu-Install-Debians.html}.

After installing ROS~2, install MAVROS and its dependencies:

\begin{lstlisting}[language=bash]
sudo apt install ros-humble-mavros ros-humble-mavros-extras
ros2 run mavros install_geographiclib_datasets.sh
\end{lstlisting}

Add the ROS~2 setup to the Jetson's shell profile so it is sourced automatically:

\begin{lstlisting}[language=bash]
echo "source /opt/ros/humble/setup.bash" >> ~/.bashrc
source ~/.bashrc
\end{lstlisting}

\textbf{Grant UART port permissions for the Jetson--Pixhawk link:}

\begin{lstlisting}[language=bash]
# Verify the UART device is present
ls /dev/ttyTHS*
# Add user to dialout group for persistent access
sudo usermod -aG dialout $USER
# Log out and back in, then verify:
ls -l /dev/ttyTHS1
\end{lstlisting}

\subsection{Camera Setup (Arducam IMX219)}

The Arducam IMX219 connects to the Jetson CSI0 port via the 22-pin to 15-pin ribbon cable. Install the Arducam driver:

\begin{lstlisting}[language=bash]
# Follow the official Arducam Jetson installation guide for the IMX219
# https://docs.arducam.com/Nvidia-Jetson-Camera/Native-Camera/Quick-Start-Guide/
\end{lstlisting}

Verify the camera is working with GStreamer:

\begin{lstlisting}[language=bash]
gst-launch-1.0 nvarguscamerasrc ! nvoverlaysink
\end{lstlisting}

\subsection{drone-brain Repository}
\label{sec:dronebrainrepo}

The \texttt{drone-brain} repository contains all autonomous mission scripts, pre-flight tests, and helper utilities. It runs on the Jetson Orin Nano.

\textbf{Clone and install:}

\begin{lstlisting}[language=bash]
git clone https://github.com/autonomous-forest-drone/drone-brain.git ~/drone-brain
cd ~/drone-brain
pip3 install -r requirements.txt
\end{lstlisting}

\textbf{Repository structure:}

\begin{tabular}{p{3.5cm}p{10cm}}
\toprule
\textbf{Directory / File} & \textbf{Contents} \\
\midrule
\texttt{missions/} & Autonomous flight mission scripts. Each script encodes a complete flight plan and is intended to be run in the field. \\
\texttt{tests/} & Commissioning and pre-flight test scripts. See Section~\ref{sec:testing} for the required test sequence. \\
\texttt{helpers/} & Utility scripts that run alongside a mission in a separate terminal. Includes telemetry monitoring and image transmission. \\
\texttt{docs/} & Additional developer documentation specific to the repository. \\
\texttt{requirements.txt} & Python dependencies. \\
\texttt{README.md} & Entry point for the repository. Always read this first -- it is the authoritative source for any setup steps not covered here. \\
\bottomrule
\end{tabular}

\textbf{All scripts follow the same two-terminal pattern:}

Terminal 1 -- start MAVROS (required for every run):
\begin{lstlisting}[language=bash]
ros2 launch mavros px4.launch fcu_url:=/dev/ttyTHS1:115200
\end{lstlisting}

Terminal 2 -- run a script:
\begin{lstlisting}[language=bash]
python3 ~/drone-brain/missions/<mission>.py
python3 ~/drone-brain/tests/<test>.py
python3 ~/drone-brain/helpers/<helper>.py
\end{lstlisting}

Use \texttt{tmux} to keep sessions alive over SSH:
\begin{lstlisting}[language=bash]
tmux new -s drone
# Split: Ctrl+B then %
# Detach: Ctrl+B then D
# Reattach: tmux attach -t drone
\end{lstlisting}

\subsection{ground-control-station Repository}

The \texttt{ground-control-station} repository runs on the monitoring laptop. It receives and displays data transmitted by the drone.

\begin{lstlisting}[language=bash]
git clone https://github.com/autonomous-forest-drone/ground-control-station.git \
    ~/ground-control-station
cd ~/ground-control-station
pip3 install -r requirements.txt
\end{lstlisting}

Refer to the repository README for the current list of available scripts. At the time of writing, it includes an \texttt{image-receiver} script that displays images transmitted by the drone's \texttt{image-sender} helper.

% =============================================================
\section{Commissioning and Testing}
\label{sec:testing}
% =============================================================

After completing hardware assembly, wiring, Pixhawk configuration, and software setup, commission the drone by working through the test sequence below in order. Do not advance to the next test until the current one passes cleanly.

\subsection{Safety Setup for Outdoor Testing}

For all initial outdoor tests, use physical restraints to limit flyaway risk:

\begin{itemize}
    \item Drive four pegs into the ground at the corners of a roughly 1\,m square centred on the drone's takeoff position.
    \item Attach a rope or strong cord from each landing leg to a peg, leaving enough slack for the drone to lift approximately 0.5--1\,m off the ground.
    \item Ensure all ropes are free of the propeller arc.
    \item Keep one operator on the RC transmitter at all times, ready to switch to Altitude/Position mode.
\end{itemize}

\subsection{Test Sequence}
\label{sec:testseq}

\subsubsection{Test 1 -- Motor Test (No Propellers)}
\label{sec:motortest}

With propellers removed and the drone connected to QGC via USB:

\begin{enumerate}
    \item Go to \textbf{Vehicle Setup $\rightarrow$ Actuators $\rightarrow$ Motors}.
    \item Spin each motor individually at low throttle using the sliders.
    \item Verify each motor spins in the correct direction (see motor diagram in the User Manual, Appendix~B): M1 Front-Right CCW, M2 Back-Left CW, M3 Front-Left CW, M4 Back-Right CCW.
    \item If a motor spins in the wrong direction, swap any two of its three phase wires at the ESC output.
\end{enumerate}

\subsubsection{Test 2 -- RC Channel Verification}

With the RC transmitter powered on and QGC connected:

\begin{enumerate}
    \item Go to \textbf{Vehicle Setup $\rightarrow$ Radio}.
    \item Move each stick and switch on the FS-i6 and confirm the corresponding channel bar moves in QGC.
    \item Verify stick directions match expected directions (e.g., throttle up = Channel 3 increases).
    \item Confirm Channel 5 (3-position switch) cycles through the three flight modes and Channel 6 (arming switch) responds.
\end{enumerate}

\subsubsection{Test 3 -- Arm-Only Script}

Install propellers. Move outdoors. Use the rope and peg safety setup.

\begin{enumerate}
    \item SSH into the Jetson and start MAVROS (Terminal 1).
    \item In Terminal 2, run the arm-only test script from the \texttt{tests/} directory.
    \item The drone should arm (motors spin at idle) and then disarm after a short delay without taking off.
    \item Verify QGC shows ARMED then DISARMED and no errors appear.
\end{enumerate}

\subsubsection{Test 4 -- Takeoff with Manual Intercept}

\begin{enumerate}
    \item Run the takeoff-only script from \texttt{tests/}.
    \item The drone lifts to a set hover altitude. \textbf{At any point}, switch the RC transmitter to Stabilise mode to cancel Offboard and take manual control.
    \item Practice the manual intercept several times to confirm the safety override works reliably before proceeding.
\end{enumerate}

\subsubsection{Test 5 -- Takeoff and Land}

\begin{enumerate}
    \item Run the takeoff-and-land script. The drone should arm, climb to hover altitude, hold position for a few seconds, descend, land, and disarm automatically.
    \item Verify there is no significant horizontal drift during hover.
\end{enumerate}

\subsubsection{Test 6 -- Takeoff, Move Forward, Land}

\begin{enumerate}
    \item Run the forward-movement script. The drone should take off, translate a short distance forward, then return and land.
    \item Verify the movement is smooth and the drone returns to approximately the launch position.
    \item This test validates that the Offboard velocity commands are working correctly.
\end{enumerate}

\subsubsection{Test 7 -- Full Movement Test}

\begin{enumerate}
    \item Run the movement-test script. This test takes off and performs controlled movements on all axes in sequence: forward, backward, left, right, yaw left, yaw right, and altitude step, then lands.
    \item Verify all movements are smooth and the drone responds correctly on each axis.
    \item This is the primary test to validate full Offboard control before GPS-based navigation.
\end{enumerate}

\subsubsection{Test 8 -- GPS Position Navigation}

\begin{enumerate}
    \item Physically carry the drone to the intended goal position. Run the goal-setting command with the \texttt{--set-goal} flag to record the GPS coordinates:
\begin{lstlisting}[language=bash]
python3 ~/drone-brain/tests/<gps-nav-test>.py --set-goal
\end{lstlisting}
    \item Return the drone to the launch position.
    \item Run the GPS navigation test (without \texttt{--set-goal}). The drone should navigate autonomously to the recorded coordinates, hold for a moment, and return to launch.
    \item Refer to the test script file and its inline documentation for the exact flags and behaviour.
\end{enumerate}

\subsubsection{Test 9 -- Camera Tests}

\begin{enumerate}
    \item Run the camera test scripts from \texttt{tests/} to verify image capture. These can be run on the ground with the drone stationary.
    \item Check that captured frames are correctly saved or transmitted.
    \item Run the telemetry monitoring helper and confirm live telemetry messages are printing to the terminal and/or appearing on the ground station.
\end{enumerate}

% =============================================================
\section{Known Limitations}
% =============================================================

\subsection{GPS Accuracy Under Forest Canopy}

The Holybro M10 GPS has a position accuracy of approximately 2.5\,m under good open-sky conditions. Under forest canopy, satellite reception degrades significantly:

\begin{itemize}
    \item The GPS position becomes \textbf{noisy and unreliable} under dense tree cover. Position errors can exceed several metres and fluctuate rapidly.
    \item The PX4 EKF fuses GPS with the barometer and IMU, but sustained GPS noise causes position hold drift in Offboard mode.
    \item Waypoint-based navigation in forest conditions therefore produces imprecise results and is not suitable for close-range obstacle avoidance.
\end{itemize}

This is a fundamental hardware limitation of single-frequency GNSS in obstructed environments, not a software issue. Addressing it would require a different positioning approach (see Section~\ref{sec:futurework}).

\subsection{Weather Dependency}

The drone has no weatherproofing. Rain, fog, or high humidity can damage electronics and corrupt sensor readings. Maximum sustained wind for safe autonomous operation is 3\,m/s.

% =============================================================
\section{Future Work}
\label{sec:futurework}
% =============================================================

The following directions are the most relevant for a developer continuing this project.

\subsection{Positioning Under Canopy}

Replacing or augmenting the GPS with sensors that work under tree cover is the most impactful improvement available:

\begin{itemize}
    \item \textbf{Optical flow sensor} (e.g.\ PX4FLOW or equivalent) -- measures ground-relative velocity from a downward-facing camera. Fuses with the barometer in PX4 EKF to provide position hold without GPS. Works indoors and under canopy but drifts over time.
    \item \textbf{Depth camera or LiDAR} (e.g.\ Intel RealSense, RPLiDAR) -- enables local obstacle detection and SLAM-based localisation on the Jetson.
    \item \textbf{RTK GPS} -- provides centimetre-level accuracy outdoors but does not solve the canopy occlusion problem.
\end{itemize}

\subsection{Image Transmission}

The current \texttt{image-sender} helper transmits 128$\times$128 pixel images over MAVLink Tunnel messages every 20 seconds. Image quality and update rate are limited by the MAVLink channel bandwidth. Alternatives include:
\begin{itemize}
    \item Transmitting over a dedicated WiFi link (requires the drone to remain within WiFi range).
    \item Logging full-resolution images to the Jetson's microSD and retrieving them post-flight.
\end{itemize}

% =============================================================
\section{Repository and Access Summary}
% =============================================================

\begin{tabular}{p{4.5cm}p{9cm}}
\toprule
\textbf{Resource} & \textbf{Location} \\
\midrule
GitHub organisation & \url{https://github.com/autonomous-forest-drone} \\
drone-brain (Jetson scripts) & \url{https://github.com/autonomous-forest-drone/drone-brain} \\
ground-control-station & \url{https://github.com/autonomous-forest-drone/ground-control-station} \\
User Manual source & \url{https://github.com/autonomous-forest-drone/user-manual} \\
PX4 parameter file & \texttt{drone-brain} repo root or dedicated config repo in the organisation \\
3D print files & Project GitHub organisation \\
\bottomrule
\end{tabular}

\begin{thebibliography}{9}
\bibitem{px4}
    PX4 Development Team,
    \emph{PX4 Autopilot User Guide},
    \url{https://docs.px4.io}, 2024.
\bibitem{ros2humble}
    Open Robotics,
    \emph{ROS 2 Humble Hawksbill Documentation},
    \url{https://docs.ros.org/en/humble/}, 2024.
\bibitem{mavros}
    MAVROS Contributors,
    \emph{MAVROS -- MAVLink extendable communication node for ROS 2},
    \url{https://github.com/mavlink/mavros}, 2024.
\bibitem{holybro_s500}
    Holybro,
    \emph{S500 V2 Development Kit User Guide},
    \url{https://holybro.com/products/s500-v2-development-kit}, 2024.
\bibitem{arducam_imx219}
    Arducam,
    \emph{IMX219 Camera for NVIDIA Jetson -- Quick Start Guide},
    \url{https://docs.arducam.com/Nvidia-Jetson-Camera/Native-Camera/Quick-Start-Guide/}, 2024.
\end{thebibliography}

\end{document}
